\documentclass{article}

\usepackage{../../Style}

\title{Algebraic Geometry}
\author{Shing Tak Lam}
\date{Michaelmas 2023\footnote{Based on lectures by Dhruv Ranganathan. Last updated \today.}}

\makeindex

\begin{document}

\maketitle

\tableofcontents

\lec

\subsection{Plan}

\begin{enumerate}
    \item Basics of sheaves of topological spaces
    \item Definitions of schemes and of morphisms
    \item Properties of schemes (the algebraic geometry analogues of compactness)
    \item Rapid introduction to the cohomology of sheaves
\end{enumerate}

\subsection{Resources}

\begin{enumerate}
    \item Course webpage - Google Dhruv
    \item Books:
        \begin{itemize}
            \item Hartshorne - Algebraic Geometry, Chapter 2 and \(\epsilon\times\)Chapter 3
            \item Vakil - Rising Sea, good online (free) alternative
            \item Eisenbud-Harris - Geometry of schemes
        \end{itemize}
    \item Four examples sheets on webpage
    \item SAGES - Look on Dhruv's webpage
\end{enumerate}

\subsection{Why scheme theory?}

\textbf{Moduli theory} -- It is better to study families of varieties than one at a time. Or better, all varieties of a given type at the same time.

\underline{Examples of moduli:} The set of all lines in \(\P^2\). A line in \(\bb P^2\) is
\[\{a X_0 + bX_1 + cX_2 = 0\}\]
In particular, as lines are paramterised by the triple \((a,b, c)\), we have a correspondence
\[\P^2_\text{dual} = \text{lines in }\P^2 \bij \bb P^2\]
The same logic applies for degree \(d\) hypersurface in \(\P^n\), and we get
\[\text{degree }d\text{ hypersufaces in }\P^n \bij \P^N\]
where \(N = {n + d \choose d} - 1\). However, there is something wrong with this picture. Some polynomials are of the form
\[f = f_1^2f_2\]
But in this case,
\[\van(f) = \van(f_1f_2)\]
and \(f_1f_2\) is (in general) not of degree \(d\).

\underline{A solution} Take \(U_d \sse \P^N\), where
\[[f] \in U_d \iff f\text{ has no repeated factors}\]
But \(U_d\) is not compact.

\underline{Output of scheme theory} for a fixed projective space \(\P^n\), we obtain a space
\[\text{subvarieties of }\P^n \bij \Var(\P^n) \subsetneq \Hilb(\P^n) \bij \text{ subschemes of }\P^n\]
where \(\Hilb(\P^n)\) is the Hilbert scheme of \(\P^n\), which is compact in the Euclidean topology. That is, the limit of varieties need not be a variety, but limits of schemes are always schemes.

In scheme theory,
\[\van(X_0 + X_1 + X_2) \text{ and }\van((X_0 + X_1 + X_2)^2)\]
are not isomorphic as schemes.

\textbf{Weil conjectures}\index{Weil conjectures} Fix \(f \in \Z[X_0, \dots, X_{n+1}]\) homogeneous. We have
\[X = \van(f) \sse \P^{n+1}_\C\]
We will assume that \(X\) is smooth. That is, \(X\) is a (complex) manifold. In particular, \(X\) is a compact topological space, and so we have numbers \(b_0(X), \dots, b_{2n}(X)\) called the Betti numbers, where
\[b_i(X) = \rank(H_i(X, \Z))\]
In particular, we have the Euler characteristic \(\chi(X)\).

On the other hand, fix a prime number \(p\), and let
\[N_m = \text{number of solutions to }f\text{ in }\bb F_{p^m} = \abs{X(\bb F_{p^m})}\]
Define the Weil \(\zeta\) function
\[\zeta(X; t) = \exp(\sum_m\frac{N_m}{m}t^m)\]

\begin{theorem}
    [Weil conjectures, Grothendieck]
    \begin{enumerate}
        \item \(\zeta(X; t)\) is a rational function in \(t\)
        \item Moreover,
        \[\zeta(X; t) = \frac{P_0(t)P_2(t)\cdots P_{2n}(t)}{P_1(t)P_3(t)\cdots P_{2n-1}(t)}\]
        where \(\deg(P_i(t)) = b_i(X)\).
    \end{enumerate}
\end{theorem}

\underline{Upshot:} The proof is (fundamentally) via scheme theory. We need a space \(\mf X\) which interpolates geometry over \(\C\) and geometry over finite fields.

\printindex

\end{document}